% =============================================================================
%  Literature Review Chapter — [Thematic Cluster Name]
%
%  This file is a template for one thematic literature chapter.
%  Duplicate it for each of your 4–5 literature clusters (e.g., sections 4–8).
%  Rename the file and update the \section label and all \label{} values.
%
%  Hints
%  -----
%  - Chapter intro (2–3 sentences): state what this chapter covers, which
%    corpus groups it draws from, and how it relates to the research question.
%  - Each subsection = one conceptual sub-topic or methodological family.
%    Do NOT organise by paper — organise by concept or technique.
%  - Each subsection should read as critical synthesis, not annotation.
%    Ask: what does this group of papers collectively establish or fail to establish?
%  - Final subsection ("Positioning against the [target setting]"): this is
%    mandatory. Explicitly state what this body of work contributes to your
%    research direction and what it leaves open. This feeds directly into the
%    gap analysis in Section~\ref{sec:synthesis_gap}.
%  - Aim for 6–10 subsections per chapter.
%  - Typical length: 4–8 pages per chapter.
% =============================================================================

\section{[Thematic Cluster Name — e.g., Adapter Composition and Reuse]}
\label{sec:lit_cluster_N}
% Rename the label to match your cluster number, e.g.:
%   sec:lit_cluster_1, sec:lit_cluster_2, ... sec:lit_cluster_5

This chapter [REPLACE: briefly describe what this chapter covers, which corpus
groups it draws from (e.g., G2), and how it connects to the main research question].

% =============================================================================
\subsection{[Conceptual Sub-Topic 1 — e.g., From Single-Task to Multi-Task Methods]}
\label{sec:lit_cluster_N:subtopic1}

% Hints:
% - Open with the structural observation that motivates the sub-topic.
% - Group papers by what they share conceptually, not chronologically.
% - Name each paper with \citet{} and follow immediately with its contribution.
% - Identify what property or limitation is shared across this group.

[REPLACE: Describe the conceptual starting point for this sub-topic.
Explain why the PEFT / system / protocol literature produced this particular
line of work.]

\citet{ExamplePaper2020} introduced [method name], which [describe the key
contribution in one sentence]. \citet{ExamplePaper2021} extended this by
[describe the extension]. Together, these works establish that [state the
collective finding or implication for your research direction].

% =============================================================================
\subsection{[Conceptual Sub-Topic 2]}
\label{sec:lit_cluster_N:subtopic2}

% Hints:
% - If the sub-topic involves a formal mechanism, include a named equation.
% - If the sub-topic introduces a taxonomy, use \itemize or \enumerate.

[REPLACE: Introduce the second sub-topic and its key papers.]

% --- Example: inline equation for a key mechanism ---
% Replace or delete as needed.
\begin{equation}
  h' = \mathrm{Attention}(h \cdot W^Q,\; [h_1,\ldots,h_K] \cdot W^K,\;
                           [h_1,\ldots,h_K] \cdot W^V)
  \label{eq:cluster_N_example}
\end{equation}

[REPLACE: Explain what the equation represents and what structural property it
captures. Note any limitations relevant to your research direction.]

% =============================================================================
\subsection{[Conceptual Sub-Topic 3]}
\label{sec:lit_cluster_N:subtopic3}

[REPLACE: Continue developing the literature narrative. Where sub-topics follow
a logical progression (e.g., first X, then X+Y, then X+Y+Z), make that
progression explicit with transition sentences.]

% =============================================================================
\subsection{[Conceptual Sub-Topic 4]}
\label{sec:lit_cluster_N:subtopic4}

[REPLACE: Continue as above.]

% =============================================================================
% Optional: add further subsections as needed.
% \subsection{[Conceptual Sub-Topic 5]}
% \label{sec:lit_cluster_N:subtopic5}

% =============================================================================
\subsection{Positioning against the [Target Setting / Research Direction]}
\label{sec:lit_cluster_N:positioning}

% Hints:
% - This is the most important subsection in each literature chapter.
% - DO NOT summarise what you just said. Instead, ask: what does this whole
%   body of work contribute toward the research direction, and what does it
%   presuppose or leave unaddressed?
% - Identify the gap clearly. Use language like:
%     "No reviewed work addresses X."
%     "Every system reviewed here assumes Y, which is incompatible with Z."
%     "The obstacle is not [trivial sub-problem] — it is [the real structural problem]."
% - This text feeds directly into Section~\ref{sec:synthesis_gap}.

[REPLACE: Describe what this body of work contributes toward your research
direction. Then identify the key limitation or assumption that prevents it from
fully addressing your research question.]

The [methods / systems / protocols] reviewed in this chapter [REPLACE: describe
the collective contribution]. However, [REPLACE: identify the structural
limitation — e.g., "every reviewed system assumes a centralised X" or "no
reviewed method addresses Y in the absence of Z"]. This limitation defines
[REPLACE: name the gap], which is examined in the synthesis in
Section~\ref{sec:synthesis_gap}.
